\chapter{La familia, comunidad de amor}

\epigraph{«Donde hay dos o tres reunidos en mi nombre, allí estoy yo en medio de ellos.»}{Mateo 18,20}

Este libro empezó con un bebé.

Un bebé que nacía con la cabeza sin cerrar, incapaz de sostenerla por sí mismo, incapaz de sobrevivir una sola noche sin que alguien le abrazara, le alimentara, le mantuviera con vida. Dijimos entonces que esa fragilidad no era un fallo de diseño, sino el diseño mismo. Que nacemos inacabados porque estamos hechos para ser completados por otros. Que el amor no es un adorno de la vida humana, sino su condición de posibilidad.

Ahora, al final del camino, quiero que volvamos a ese bebé. Pero esta vez vamos a ampliar el encuadre. Vamos a ver la foto completa.

Porque ese bebé no estaba solo con su madre. Nunca lo estuvo.

\section{La tribu que sostiene al niño}

Hay un proverbio africano que se ha hecho famoso ---tanto que ya nadie sabe a qué pueblo pertenece exactamente---: «Hace falta una aldea entera para criar a un niño.» Como muchos proverbios, dice en una línea lo que la ciencia ha necesitado décadas para demostrar.

En el capítulo 1 hablamos del dilema obstétrico: la pelvis se estrecha, el cerebro crece, el bebé nace inmaduro. La madre no puede sola. Necesita ayuda. Y dijimos que el padre que se queda, el que cuida, el que baja su testosterona y sube su oxitocina, fue una pieza clave en la supervivencia de nuestra especie.

Pero la historia no termina ahí. Hay otra pieza del rompecabezas que entonces solo mencionamos de pasada y que ahora merece toda nuestra atención.

Se llama la \textbf{hipótesis de la abuela}.\footnote{La hipótesis fue propuesta formalmente por la antropóloga Kristen Hawkes, de la Universidad de Utah, a partir de sus estudios con los Hadza de Tanzania en los años noventa. Cf. Hawkes et al., ``Grandmothering, menopause, and the evolution of human life histories'', \textit{PNAS}, 1998.}

El enigma es el siguiente: la menopausia humana es una rareza biológica. En casi todas las especies, las hembras son fértiles hasta el final de su vida. ¿Por qué las mujeres humanas dejan de ser fértiles décadas antes de morir? Desde un punto de vista puramente evolutivo, ¿no sería más «eficiente» seguir reproduciéndose?

La respuesta que propone esta hipótesis es fascinante: las abuelas que dejaban de tener hijos propios y se dedicaban a ayudar a criar a los hijos de sus hijas contribuían más a la supervivencia de sus genes que si hubieran seguido teniendo crías. Una abuela que recolecta alimentos, que vigila a los niños, que cuida a la madre agotada, que enseña a la nieta a identificar raíces comestibles, que consuela al nieto asustado, multiplica las posibilidades de que esos nietos sobrevivan y se reproduzcan.

La evolución, en otras palabras, no favoreció solo a la madre que cuidaba. Ni solo a la pareja que permanecía unida. Favoreció a la \emph{familia extensa}. A la red de cuidado intergeneracional. A la tribu.

\section{Las abuelas de Dios}

Hay algo que me conmueve profundamente de esta hipótesis, y es que la ciencia, sin pretenderlo, está describiendo algo que la fe ya conocía.

Porque lo que la hipótesis de la abuela dice, traducido al lenguaje de la fe, es que el ser humano está diseñado para que el amor desborde. No basta el amor de la madre. No basta el amor de la pareja. Hace falta un amor que se extienda más allá, que cruce generaciones, que cree redes donde cada persona cuida de otra y es cuidada a su vez.

La abuela que deja de engendrar para dedicarse a cuidar está haciendo algo que se parece mucho a lo que la tradición cristiana llama \emph{fecundidad espiritual}: una vida que da fruto no a través de la reproducción, sino a través de la entrega. Una vida que se multiplica cuidando.

Pensad en vuestras propias abuelas. O en esas figuras que, sin ser abuelas de sangre, ejercieron ese papel: la vecina que os recogía del colegio, la tía que os enseñó a cocinar, el abuelo que os contaba historias que todavía recordáis. Cada una de esas personas fue, a su manera, parte de la aldea que os crió. Cada una de ellas hizo posible que llegarais hasta aquí.

La familia no es la pareja encerrada en su piso. La familia es una red. Una red que se teje con hilos de sangre, pero también de cercanía, de costumbre, de cariño. Una red que, cuando funciona, sostiene a los débiles, acoge a los nuevos, acompaña a los que se van. Y cuando no funciona ---porque a veces no funciona---, su ausencia se nota como se nota el frío cuando falta el techo.

\section{La Iglesia doméstica}

El Concilio Vaticano II, en uno de esos hallazgos de lenguaje que parecen sencillos pero encierran siglos de reflexión, llamó a la familia «Iglesia doméstica».\footnote{Concilio Vaticano II, Constitución dogmática \textit{Lumen Gentium}, n. 11 (1964): «En esta especie de Iglesia doméstica, los padres deben ser para sus hijos los primeros predicadores de la fe, mediante la palabra y el ejemplo.» Juan Pablo II retomó y desarrolló ampliamente esta idea en la Exhortación apostólica \textit{Familiaris Consortio} (1981), especialmente en los nn. 49-62.}

\textit{Iglesia doméstica.} La expresión es más fuerte de lo que parece. No dice que la familia sea «como» una iglesia, en sentido metafórico. Dice que la familia \emph{es} Iglesia. La versión más pequeña, más cotidiana, más real de la Iglesia. El lugar donde la fe no se explica, sino que se vive. Donde no se predica con palabras, sino con gestos.

Porque, pensémoslo bien: ¿dónde aprende un niño lo que es el perdón? No en una clase de catequesis. Lo aprende cuando ve a su padre pedirle perdón a su madre después de una discusión. ¿Dónde descubre lo que es la generosidad? Cuando ve a sus padres acoger en casa a alguien que lo necesita. ¿Dónde intuye, por primera vez, que existe algo más grande que él, algo que sostiene la vida y le da sentido? Quizá en esa oración sencilla antes de cenar que nadie le obliga a decir, pero que todo el mundo dice.

La familia es el primer lugar donde se experimenta ---o no se experimenta--- que el amor es real. Y esa experiencia marca para siempre. No porque determine fatalmente el futuro, sino porque establece el suelo sobre el que todo lo demás se construye. Un niño que ha visto amor sabe que el amor existe, aunque luego le cueste encontrarlo. Un niño que no lo ha visto tendrá que descubrirlo por sí mismo, que es más difícil ---pero no imposible, porque la gracia llega donde la naturaleza no alcanza---.

\section{Los suegros y otras pruebas de fe}

Y ahora viene el momento de hablar de algo que todo el que se ha casado conoce, que todos los manuales de pastoral mencionan con delicadeza y que la vida real sirve sin anestesia: las familias de origen.

Porque cuando dos personas se casan, no se casan solas. Se casan con sus historias, sus costumbres, sus formas de celebrar la Navidad y de doblar las toallas. Se casan con sus padres, sus hermanos, sus tradiciones familiares. Y todo eso tiene que encontrar un encaje que, a veces, rechina.

El Génesis lo dice con una claridad que impresiona: «Dejará el hombre a su padre y a su madre y se unirá a su mujer.»\footnote{Gn 2,24.} Dejar. No abandonar, no renegar, no olvidar. Dejar. Soltar el lugar de hijo para ocupar el lugar de esposo. Soltar el lugar de hija para ocupar el lugar de esposa. Crear un espacio nuevo, una casa nueva ---aunque sea el mismo piso de alquiler---, un «nosotros» que es distinto del «nosotros» de la familia en la que creciste.

Y eso, que suena muy bonito en un libro, en la práctica significa aprender a navegar entre dos familias que os quieren, que quieren lo mejor para vosotros, que a veces no se ponen de acuerdo sobre qué es lo mejor para vosotros, y que ---seamos honestos--- a veces meten la nariz donde no las llaman.

No voy a dar recetas, porque cada familia es un mundo y cada suegra es un universo.\footnote{Y cada suegro. Los suegros también existen, aunque la cultura popular haya concentrado todo su humor en las suegras. Quede aquí constancia de que el desafío es ecuménico.} Pero sí voy a decir dos cosas que he visto confirmadas una y otra vez.

La primera: \textbf{honrar a vuestros padres no significa obedecerlos en todo}. Sois adultos. Habéis fundado una familia nueva. Las decisiones sobre vuestra casa, vuestra economía, vuestra forma de educar a vuestros hijos, son \emph{vuestras}. Pedir consejo es sabio. Dejarse gobernar es un error. Y cuando haya que poner un límite, es mejor que lo ponga el hijo o la hija de sangre, no el cónyuge. Si tu madre se mete demasiado, eres tú quien tiene que hablar con ella. No tu marido. Eso no es cobardía: es cuidar a las dos personas que quieres.

La segunda: \textbf{la familia política es un regalo que no habéis elegido, pero que podéis aprender a querer}. No os casasteis con ellos, pero forman parte de la vida de la persona con la que sí os casasteis. Y esa persona es, en parte, lo que es gracias a ellos. Así que miradlos con agradecimiento, aunque a veces os cueste. Y tened paciencia. Mucha paciencia. La misma que esperáis que ellos tengan con vosotros.

\section{La puerta abierta}

Si hay algo que define a una familia cristiana, no es que rece mucho ni que vaya a misa todos los domingos ---aunque ambas cosas estén muy bien---. Lo que define a una familia cristiana es que su puerta está abierta.

Abierta al hijo que llega, esperado o por sorpresa. Abierta al amigo que necesita un plato de comida y una conversación. Abierta al vecino que está solo. Abierta al sobrino que necesita un lugar donde quedarse mientras encuentra su camino. Abierta al mundo.

Hablamos en el capítulo 10 de la fecundidad como irradiación: un matrimonio fecundo es el que genera vida a su alrededor.\footnote{Cf. capítulo 10 de este libro.} Pues la familia es el lugar donde esa fecundidad se hace concreta. Una familia abierta es una familia que no vive para sí misma. Que no se agota en sus propias necesidades. Que tiene la generosidad de mirar más allá de sus paredes y preguntarse: ¿a quién podemos acoger? ¿A quién podemos servir?

No hace falta ser héroes. A veces basta con invitar a cenar a esa pareja que acaba de mudarse al barrio y no conoce a nadie. Con acompañar al anciano del tercero al médico. Con preparar la habitación de los hijos ---que ya se fueron--- para quien la necesite. La hospitalidad no es un gran gesto: es una disposición del corazón que se traduce en cosas muy pequeñas.

San Juan Pablo II escribió que la familia está llamada a ser «la primera y vital célula de la sociedad».\footnote{Juan Pablo II, Exhortación apostólica \textit{Familiaris Consortio}, n. 42 (1981), citando el Decreto \textit{Apostolicam Actuositatem} del Concilio Vaticano II, n. 11.} No la única célula, pero sí la primera. El lugar donde se aprende a convivir, a compartir, a resolver conflictos sin destruirse. Si la familia funciona, la sociedad tiene una oportunidad. Si la familia se cierra sobre sí misma, se asfixia. Como la semilla que no rompe la tierra: tiene toda la vida dentro, pero si no se abre, no da fruto.

\section{El perdón de cada día}

Hay una frase que, si la familia fuera una empresa, sería su lema corporativo. No es «te quiero», aunque esa es importante. No es «gracias», aunque esa es imprescindible. La frase más necesaria en una familia, la que más se usa en las casas donde la gente es feliz de verdad, es:

\textbf{«Perdona.»}

No hablo del gran perdón, el perdón dramático, el perdón de las películas donde alguien ha hecho algo terrible y el otro, tras una noche de lluvia y violines, decide perdonar. Ese perdón existe, y a veces hace falta. Pero no es el que sostiene una familia.

El que sostiene una familia es el perdón pequeño. El de cada día. El «perdona, he sido brusco». El «perdona, no te estaba escuchando». El «perdona, me he pasado con lo que he dicho». El «perdona, tenías razón y yo no quería verlo». Ese perdón que no sale en las películas porque no es espectacular, pero que es el aceite que mantiene la maquinaria en marcha.

Vivir juntos es rozarse. Es inevitable. Compartir espacio, tiempo, dinero, decisiones, cansancio, hijos, preocupaciones ---todo eso genera fricción---. Y la fricción, si no se lubrica con perdón, acaba desgastando. No de golpe. Poco a poco. Como el agua que erosiona la piedra: no por su fuerza, sino por su constancia.

El perdón en la familia no es un acto heroico. Es una necesidad higiénica. Tan básica como ducharse. Y, como la ducha, hay que practicarlo todos los días, no una vez al año.\footnote{San Pablo lo dijo con menos humor pero más profundidad: «Que no se ponga el sol sobre vuestra ira» (Ef 4,26). En otras palabras: no dejéis que un enfado pernocte en vuestra casa. Resolvedlo antes de dormir. No siempre es posible, pero como ideal es imbatible.}

Lo contrario del perdón no es el rencor espectacular. Es la cuenta que se va acumulando en silencio. Es el «no digo nada, pero me acuerdo». Es la lista mental de agravios que un día, sin que nadie lo espere, estalla. Las familias no se rompen por una gran traición. Se rompen por mil pequeñas heridas que nadie curó.

Perdonar no es sentir que no ha pasado nada. Perdonar es decidir que lo que ha pasado no va a definir la relación. Es elegir al otro por encima del agravio. Y eso, en una familia, hay que hacerlo todos los días. A veces varias veces al día.

Es agotador. Es necesario. Es, probablemente, lo más parecido a amar como Dios ama que podemos hacer los seres humanos corrientes.

\section{La familia imperfecta}

Y aquí quiero decir algo que ojalá alguien me hubiera dicho a mí antes de casarme: \textbf{vuestra familia no va a ser perfecta. Y eso no es un fracaso. Es lo normal.}

Habrá días en los que todo fluya: la casa estará recogida, los niños serán adorables, la cena saldrá rica y os miraréis a los ojos y pensaréis: «Qué suerte tengo.» Pero habrá muchos más días en los que la casa será un caos, los niños se pelearán, la cena se quemará y os miraréis a los ojos y pensaréis: «¿Cómo hemos llegado hasta aquí?»

Ambos días son reales. Ambos son vuestra familia.

La tentación más peligrosa no es la infidelidad ni el desamor. La tentación más peligrosa es la comparación. Es mirar a esa otra familia ---la del amigo, la del vecino, la que sale en Instagram--- y pensar que ellos sí lo tienen todo resuelto. Que su casa siempre está limpia, sus hijos siempre obedecen, su matrimonio siempre brilla. Es mentira, por supuesto. Nadie publica las fotos de la discusión de las once de la noche ni del niño llorando en el supermercado. Pero la mentira es eficaz, y siembra la duda: «Quizá lo estamos haciendo mal.»

No. No lo estáis haciendo mal. Lo estáis haciendo como se hace: con lo que tenéis, con lo que sabéis, con lo que podéis. Con más buena voluntad que aciertos, y eso es suficiente. Porque la familia no es un escaparate. Es un taller. Un lugar donde se trabaja, donde se falla, donde se vuelve a intentar. Donde la materia prima es el amor imperfecto de personas imperfectas que, sin embargo, han decidido estar juntas.

El papa Francisco, con esa franqueza suya que a veces incomoda y siempre ilumina, lo dijo así: «No existe la familia perfecta. No hay maridos perfectos, ni esposas perfectas. No nos hablemos de suegras perfectas. [\ldots] Solo existen familias en camino.»\footnote{Papa Francisco, audiencia general del 16 de septiembre de 2015.} En camino. No en destino. La perfección no es el punto de partida ni el requisito. Es el horizonte hacia el que se camina, sabiendo que nunca se llega del todo, y que eso está bien.

\section{Diseñados para amar}

Permitidme que, al cerrar este capítulo ---y con él, este libro---, vuelva una vez más al principio.

Empezamos hablando de un bebé que nacía indefenso, con la cabeza sin soldar y los ojos que apenas veían. Dijimos que esa fragilidad era el origen de todo: de la necesidad de cuidar, de vincularse, de amar. La biología nos mostraba un cuerpo diseñado para la relación: hormonas que empujan al vínculo, cerebros que se transforman con la paternidad y la maternidad, una química que favorece la permanencia.

Después la fe nos dio la clave de lectura: no estamos hechos para amar por casualidad ni por presión evolutiva. Estamos hechos para amar porque quien nos hizo es Amor. Somos imagen de un Dios que es, en sí mismo, comunión, entrega, don. Y amando ---entregándonos, saliéndonos de nosotros mismos, haciéndonos don para el otro--- nos parecemos a Él.

Vimos que el matrimonio no es un contrato con fecha de caducidad, sino una alianza sin letra pequeña. Que la indisolubilidad no es una cadena, sino la promesa más grande que una persona puede recibir. Que la fidelidad es fe, la fe absoluta que el otro deposita en ti. Que la fecundidad desborda los hijos y se derrama en todo lo que toca una familia que vive abierta. Que la sexualidad es un lenguaje que dice con el cuerpo lo que la alianza dice con palabras: «Todo lo mío es tuyo, sin reservas, para siempre.»

Y ahora, al final, vemos el cuadro completo: aquel bebé indefenso crece. Es sostenido por una madre, por un padre, por unos abuelos, por una comunidad que le cuida, le enseña, le quiere. Poco a poco aprende a hablar, a caminar, a pensar. Aprende ---si tiene suerte--- lo que es ser querido. Y un día, ya adulto, mira a otra persona a los ojos y siente algo que reconoce sin que nadie se lo explique: el deseo de entregarse. La vocación al amor. Y dice: «Sí, quiero.»

Y empieza de nuevo. Un hogar. Una alianza. Una vida compartida. Y quizá, un día, otro bebé indefenso que necesitará una cabeza que le sostengan, unos brazos que le abracen, una aldea entera que le críe. El ciclo continúa. La misma historia, en otra clave, con otros rostros, con los mismos miedos y la misma esperanza.

\medskip

Dijimos al hablar del matrimonio que la Cruz no es principalmente sufrimiento: es entrega.\footnote{Cf. capítulo 7 de este libro.} Que la persona a la que amas es tu cruz ---la razón por la que entregas tu vida--- y que Dios nos da una cruz a la que podemos abrazar cada noche. Pues bien: la familia es la extensión de ese abrazo. Es la Cruz que se hace mesa, que se hace casa, que se hace sobremesa un domingo, que se hace discusión un lunes y reconciliación un martes. Es la entrega hecha cotidianidad.

No será perfecta. Será real. Habrá días luminosos y noches oscuras. Habrá risas y portazos, cenas memorables y leche derramada, oraciones a medias y silencios que pesan. Habrá momentos en los que pensaréis que no podéis más, y momentos en los que sabréis, con una certeza que no necesita argumentos, que esto ---exactamente esto, con todo su desorden--- es lo más importante que habéis hecho en la vida.

Porque estáis diseñados para esto. No para la perfección, sino para el amor. No para acertar siempre, sino para volver a intentarlo cada mañana. No para ser una familia de postal, sino para ser una familia de verdad: torpe, ruidosa, desordenada, y llena ---llena hasta los bordes--- de algo que se parece mucho a Dios.

\medskip

Y quiero dejaros con una última imagen. Pensad en el camino de una vida entera. Desde aquel bebé con la cabeza sin cerrar hasta el último día. Es un camino largo, lleno de etapas, de giros, de subidas y bajadas. Pues bien: vuestro cónyuge es la persona que camina con vosotros ese camino. No un tramo. No una etapa. El camino entero. Hasta la misma puerta del cielo.

Porque si lo que dijimos en el capítulo 7 es verdad ---que tu esposo, tu esposa, es tu cruz, y que la Cruz es el camino hacia el Padre---, entonces tu cónyuge es, literalmente, quien te lleva al cielo. No en sentido figurado. En el sentido más real y más hondo que existe. Cada día que os amáis, cada perdón que os dais, cada sacrificio silencioso, cada noche en la que elegís al otro aunque cueste, es un paso más en ese camino de santidad que recorréis juntos. Tu cónyuge te acompaña, te sostiene, te empuja, te corrige, te abraza. Y al final del camino ---cuando llegue el momento de cruzar esa puerta definitiva---, será quien te tome de la mano y te la ponga en la mano de Dios para empezar la vida nueva.

Del bebé indefenso del primer capítulo al esposo que llega a las puertas del cielo de la mano de quien amó: esa es la historia completa. Esa es la historia para la que estáis diseñados.

Que así sea.

\paracaminar

\begin{enumerate}
  \item Pensad juntos en vuestra «aldea»: ¿quiénes son las personas ---de vuestra familia o de fuera de ella--- que os sostienen, que os cuidan, que hacen posible vuestra vida en común? ¿Les habéis dado las gracias alguna vez?

  \item El perdón diario: ¿cómo gestionáis los roces y los pequeños conflictos en vuestra vida juntos? ¿Hay alguna cuenta pendiente, alguna herida pequeña que no se ha curado, que podríais atender hoy? Antes de cerrar este libro, miraos a los ojos y decid en voz alta algo que necesitéis perdonar o algo por lo que necesitéis pedir perdón.

  \item Este libro se titula \textit{Diseñados para amar}. Ahora que lo habéis recorrido juntos, escribid cada uno en un papel ---sin enseñárselo al otro--- qué ha cambiado en vuestra forma de ver el matrimonio, la familia o el amor. Después, intercambiad los papeles y leed en silencio lo que ha escrito el otro. No hace falta comentar nada. A veces basta con saber.
\end{enumerate}

\vinetacapitulo{ch12}
